\chapter{Una misma frontera: flujo interior y tensión gauge}
\label{chap:hoja-comun-ns-ym}
\index{Navier--Stokes}\index{Yang--Mills}
\index{operador de frontera}\index{Stokes}

La hidrodinámica y Yang--Mills no se reúnen comparando desde fuera dos
ecuaciones célebres. Proceden de dos representaciones del mismo operador de
frontera. La orientación que el lector interior interpreta como circulación
se convierte, bajo el lector reflejado, en holonomía; la memoria simétrica
que controla el flujo se convierte en kernel positivo de transferencia.

\section{El estado común de frontera}

Sea \(K_m\) un complejo finito orientado, \(B_m\) su álgebra de datos de
frontera y \(d_m\) su coborde. El estado enriquecido es
\[
 s_m=(b_m,U_m,\theta_m,c_m,\Omega_m),
\]
donde \(U_m\) transporta, \(\theta_m^2=I\) invierte la orientación,
\(c_m\) conserva memoria y \(\Omega_m\) es una \(1\)-cocadena de conexión.
Los proyectores de paridad son
\[
 E_\pm=\frac12(I\pm\theta_m).
\]
La parte \((U_m+U_m^\dagger)/2\) es simétrica; sólo es un proyector si además
se demuestra su idempotencia.

Sobre el complejo de cocadenas se construyen dos lectores lineales:
\[
 \rho_{{\rm int},m}:C^0(K_m)\longrightarrow
 C^1(K_m)\oplus C^0(K_m),
\qquad
 \rho_{{\rm int},m}(v)=(d_mv,d_m^\ast d_mv),
\]
\[
 \rho_{{\rm ref},m}:C^1(K_m)\longrightarrow
 C^0(K_m)\oplus C^1(K_m),
\qquad
 \rho_{{\rm ref},m}(a)=(d_m^\ast a,d_md_m^\ast a).
\]
El primero transporta un potencial de vértices a su flujo de aristas y a su
respuesta interior; el segundo transporta una cocadena de aristas a su
respuesta de borde y a su operador reflejado. Los codominios son distintos,
pero ambos se construyen con el mismo coborde. Circulación, vorticidad,
holonomía y curvatura aparecen después mediante la identidad de Stokes; no
se presuponen dentro de los símbolos \(\rho\).

\section{El operador mínimo y su espectro}

En el ciclo orientado de tres vértices, sea
\[
 d_0=
 \begin{pmatrix}
 -1&1&0\\
 0&-1&1\\
 1&0&-1
 \end{pmatrix}.
\]
Los Laplacianos de Hodge de vértices y aristas son
\[
 \mathcal L_{\rm H}^{(0)}=d_0^\ast d_0,
 \qquad
 \mathcal L_{\rm H}^{(1)}=d_0d_0^\ast,
\]
y un cálculo directo da
\[
 \mathcal L_{\rm H}^{(0)}=\mathcal L_{\rm H}^{(1)}=
 \begin{pmatrix}
 2&-1&-1\\
 -1&2&-1\\
 -1&-1&2
 \end{pmatrix}
 =3I-J.
\]
En esta residencia, los superíndices \((0)\) y \((1)\) designan los
Laplacianos de Hodge sobre vértices y aristas. No designan los operadores de
elevación \(L_0,L_1\) de la cadena \(w_6\to w_{30}\). Para
\(\mathcal L_{\rm cyc}^{\rm TPK}:=2I-C-C^{-1}=3I-J\), la realización del
operador cíclico del TPK se tipa mediante
\[
 \mathfrak R_{\rm hyd}(\mathcal L_{\rm cyc}^{\rm TPK})
 =\mathcal L_{\rm H}^{(0)},\qquad
 \mathfrak R_{\rm YM}(\mathcal L_{\rm cyc}^{\rm TPK})
 =\mathcal L_{\rm H}^{(1)}.
\]

\begin{theorem}[Operador común de frontera]
\label{thm:operador-comun-frontera-ns-ym}
En el ciclo de tres celdas,
\[
 \operatorname{Spec}(\mathcal L_{\rm H}^{(0)})
 =\operatorname{Spec}(\mathcal L_{\rm H}^{(1)})
 =\{0,3,3\}.
\]
Después de retirar el modo constante ---media cero en el lector interior y
calibre en el lector de conexión--- ambos sectores poseen gap exacto \(3\).
En todo complejo finito, \(d^\ast d\) y \(dd^\ast\) tienen el mismo espectro
no nulo, con multiplicidades.
\end{theorem}

\begin{proof}
El vector \((1,1,1)\) genera el núcleo de \(3I-J\). En su complemento
ortogonal, \(J=0\), de modo que el operador es \(3I\). Para el enunciado
general, una descomposición singular \(d=U\Sigma V^\ast\) produce
\[
 d^\ast d=V\Sigma^\ast\Sigma V^\ast,
 \qquad
 dd^\ast=U\Sigma\Sigma^\ast U^\ast.
\]
Los cuadrados de los valores singulares no nulos, incluidas sus
multiplicidades, coinciden.
\end{proof}

\section{Lectura hidrodinámica}

Definimos el paso implícito
\[
 S_{\rm NS}=(I+\mathcal L_{\rm H}^{(0)})^{-1}.
\]
En el ciclo mínimo,
\[
 S_{\rm NS}=\frac14(I+J).
\]

\begin{theorem}[Contracción global del flujo finito]
\label{thm:contraccion-flujo-finito}
Para todo \(v_0\) de media cero, la iteración
\(v_{k+1}=S_{\rm NS}v_k\) existe y es única para todo \(k\geq0\), y
\[
 v_k=4^{-k}v_0,
 \qquad
 \|v_k\|\leq4^{-k}\|v_0\|,
 \qquad
 \|v_k\|^2=16^{-k}\|v_0\|^2.
\]
\end{theorem}

\begin{proof}
En media cero se tiene \(Jv=0\), luego
\(S_{\rm NS}v=v/4\). La fórmula se obtiene por inducción.
\end{proof}

\Needspace{4\baselineskip}
Una no linealidad orientada \(N_m(v,c)\) conserva la energía cuando
\[
 \langle v,N_m(v,c)\rangle=0.
\]
La advección redistribuye entonces la energía y el operador positivo controla
su propagación. La memoria \(c_m\) no se absorbe en una viscosidad elegida:
forma parte del estado que debe transportarse entre escalas.

\begingroup
\setlength{\parskip}{0pt}
\titlespacing*{\section}{0pt}{8pt plus 1pt minus .5pt}{2.5pt}
\setlength{\abovedisplayskip}{.25\baselineskip}
\setlength{\belowdisplayskip}{.25\baselineskip}
\setlength{\abovedisplayshortskip}{.12\baselineskip}
\setlength{\belowdisplayshortskip}{.18\baselineskip}
\section{Lectura Yang--Mills}

Una \(1\)-cocadena \(a\) se interpreta como potencial de conexión y posee
energía cuadrática
\[
 Q_{\rm YM}(a)=\langle a,\mathcal L_{\rm H}^{(1)}a\rangle.
\]

\begin{theorem}[Coercividad y decaimiento gauge finitos]
\label{thm:coercividad-gap-gauge-finito}
En el cociente por el modo gauge constante,
\[
 Q_{\rm YM}(a)\geq3\|a\|^2.
\]
El semigrupo \(K_t=e^{-t\mathcal L_{\rm H}^{(1)}}\) satisface
\[
 \|K_ta\|\leq e^{-3t}\|a\|
\]
en el sector físico del ciclo mínimo.
\end{theorem}

\begin{proof}
Sobre el complemento del núcleo, el
\cref{thm:operador-comun-frontera-ns-ym} da
\(\mathcal L_{\rm H}^{(1)}=3I\). Las dos fórmulas
siguen por cálculo funcional.
\end{proof}

Para una conexión no abeliana,
\[
 F_\Omega=d\Omega+\Omega\smile\Omega,
\qquad
 W_\gamma(\Omega)=\operatorname{Tr}\operatorname{Hol}_\gamma(\Omega).
\]
La positividad por reflexión requiere una factorización sobre la mitad
positiva,
\[
 \langle\theta f,Kf\rangle=\|Cf\|^2\geq0;
\]
la positividad puntual de la matriz \(K\) no la sustituye.

\section[Identidad de Stokes: circulación--vorticidad y holonomía--curvatura]{Identidad discreta de Stokes y \(2\) realizaciones: circulación--vorticidad y holonomía--curvatura}

Para una \(1\)-cocadena \(\Omega\) y una \(2\)-cadena \(K\),
\[
 \sum_{e\subset\partial K}\Omega(e)
 =
 \sum_{f\in K^{(2)}}(d\Omega)(f).
\]
Las aristas interiores se cancelan por orientación. El lector
hidrodinámico interpreta el lado izquierdo como circulación y el derecho
como vorticidad acumulada. El lector gauge interpreta el primero como
holonomía linealizada y el segundo como curvatura.

\begin{theorem}[Teorema de la hoja común de frontera]
\label{thm:hoja-comun-frontera}
La descomposición \(B_m=E_-B_m\oplus E_+B_m\) posee dos representaciones
tipadas:
\[
\boxed{
\begin{array}{ccl}
E_-B_m&\longmapsto&
\text{circulación}\ /\ \text{holonomía},\\[2mm]
E_+B_m&\longmapsto&
\text{control disipativo}\ /\ \text{kernel reflejado}.
\end{array}}
\]
La igualdad de Stokes y la coincidencia del espectro no nulo pertenecen al
objeto común. La contracción hidrodinámica y el gap gauge son consecuencias
distintas en sus respectivos codominios.
\end{theorem}

\begin{proof}
Los proyectores \(E_\pm\) son complementarios porque
\(\theta_m^2=I\). La identidad de Stokes proporciona el entrelazamiento del
sector orientado; el \cref{thm:operador-comun-frontera-ns-ym} proporciona el
espectro común. Los
\cref{thm:contraccion-flujo-finito,thm:coercividad-gap-gauge-finito}
calculan las dos realizaciones.
\end{proof}

El teorema es un cierre finito HMT: no identifica una PDE con una teoría
gauge ni confunde el gap \(3\) con una masa física. Explica por construcción
por qué regularidad interior y separación de frontera son dos funciones de
la misma memoria orientada. Cualquier paso continuo posterior debe conservar
el operador, la norma y el espectro que hacen exacta esta identidad.
\endgroup
